\subsection{混合精度 + Activation Checkpointing 的协同}

在实际大规模训练中，混合精度和 activation checkpointing 通常同时使用，形成一套完整的内存优化方案：

\begin{table}[H]
\centering
\begin{tabular}{p{0.30\textwidth}p{0.20\textwidth}p{0.20\textwidth}p{0.15\textwidth}}
\toprule
\textbf{优化技术} & \textbf{内存节省} & \textbf{计算开销} & \textbf{实现难度} \\
\midrule
混合精度 (bf16) & 激活减半 & 无（反而加速） & 低 \\
Activation checkpointing & 激活减 $\sqrt{L}\times$ & +33\% 前向 & 中 \\
梯度累积 & 等效大 batch & 无额外开销 & 低 \\
模型并行 & 线性切分 & 通信开销 & 高 \\
\bottomrule
\end{tabular}
\caption{常见内存优化技术的对比}
\end{table}

